\documentclass{article}
\usepackage[margin=3cm]{geometry}
\usepackage{amsmath}
\usepackage{amssymb}
\usepackage{amsthm}

\newtheorem{theorem}{Theorem}

\newcommand{\bb}[1]{\mathbb{#1}}
\renewcommand{\vec}[1]{\boldsymbol{#1}}

\title{Transformations of Level Sets}
\author{}
\date{}

\begin{document}

\maketitle

\section{General Transformation}

\begin{theorem}
Let $f: \bb{R}^n \to \bb{R}$ define a level set $D = L_0(f) = \left\{ \vec{r} \in \bb{R}^n \;\middle|\; f(\vec{r}) = 0 \right\}$. 
If $T: \bb{R}^n \to \bb{R}^n$ is a bijective (invertible) transformation, 
then the image set $D' = T(D)$ under $T$ is given by:
\begin{equation}
  D' = \left\{ \vec{r} \in \bb{R}^n \;\middle|\; (f \circ T^{-1})(\vec{r}) = 0 \right\}
\end{equation}
\end{theorem}

\paragraph{Informal Proof}
By definition of function composition, $ (f \circ T^{-1})(\vec{r}) = f(T^{-1}(\vec{r})) $. \\
Letting $\vec{u} = T^{-1}(\vec{r})$, we have $\vec{r} = T(\vec{u})$. Substituting back gives:
\begin{align*}
  D' &= \left\{ T(\vec{u}) \;\middle|\; f(\vec{u}) = 0 \right\} \\
  D' &= T(D)
\end{align*}

\begin{proof}
\begin{align*}
  \vec{r} \in T(D) &\iff \exists \vec{u} \in D \quad \text{such that} \quad \vec{r} = T(\vec{u}) \\
  &\iff \exists \vec{u} \in \bb{R}^n \quad \text{such that} \quad f(\vec{u}) = 0 \quad \text{and} \quad \vec{u} = T^{-1}(\vec{r}) \\
  &\iff f(T^{-1}(\vec{r})) = 0 \\
  &\iff (f \circ T^{-1})(\vec{r}) = 0
\end{align*}
Therefore, $D' = T(D) = \left\{ \vec{r} \in \bb{R}^n \;\middle|\; (f \circ T^{-1})(\vec{r}) = 0 \right\}$.
\end{proof}

\section{Examples - 2D Level Sets (Level Curves)}

Let $\vec{r} = \begin{pmatrix} x \\ y \end{pmatrix}$ and $D = \left\{ \begin{pmatrix} x \\ y \end{pmatrix} \;\middle|\; f(x, y) = 0 \right\}$.

\subsection{Translation}

For a translation vector $\Delta \vec{r} = \begin{pmatrix} \Delta x \\ \Delta y \end{pmatrix}$, the translation mapping is $T_{\Delta \vec{r}}(\vec{r}) = \vec{r} + \Delta \vec{r}$.

\begin{align*}
  D_T &= \left\{ \begin{pmatrix} x \\ y \end{pmatrix} \;\middle|\; f(x - \Delta x, y - \Delta y) = 0 \right\} \\[1ex]
  \vec{u} &= \begin{pmatrix} x' \\ y' \end{pmatrix} = \begin{pmatrix} x - \Delta x \\ y - \Delta y \end{pmatrix} = \vec{r} - \Delta \vec{r} = T_{-\Delta \vec{r}}(\vec{r}) \\[1ex]
  \vec{r} &= \vec{u} + \Delta \vec{r} = T_{\Delta \vec{r}}(\vec{u}) \\[1ex]
  D_T &= \left\{ T_{\Delta \vec{r}}\left( \begin{pmatrix} x' \\ y' \end{pmatrix} \right) \;\middle|\; f(x', y') = 0 \right\} = T_{\Delta \vec{r}}(D)
\end{align*}

\subsection{Scaling}

Let $S = \begin{pmatrix} s_x & 0 \\ 0 & s_y \end{pmatrix}$ be the scaling matrix, where $L_S(\vec{r}) = S\vec{r} = \begin{pmatrix} s_x x \\ s_y y \end{pmatrix}$.

\begin{align*}
  D_S &= \left\{ \begin{pmatrix} x \\ y \end{pmatrix} \;\middle|\; f\left( \frac{x}{s_x}, \frac{y}{s_y} \right) = 0 \right\} \\[1ex]
  \vec{u} &= \begin{pmatrix} x' \\ y' \end{pmatrix} = \begin{pmatrix} x / s_x \\ y / s_y \end{pmatrix} = S^{-1}\vec{r} = L_S^{-1}(\vec{r}) \\[1ex]
  \vec{r} &= S\vec{u} = L_S(\vec{u}) \\[1ex]
  D_S &= \left\{ S \begin{pmatrix} x' \\ y' \end{pmatrix} \;\middle|\; f(x', y') = 0 \right\} = L_S(D)
\end{align*}

\subsection{Rotation}

Let $R_{\theta} = \begin{pmatrix} \cos\theta & -\sin\theta \\ \sin\theta & \cos\theta \end{pmatrix}$. The inverse transformation matrix is:
\[
  R_{\theta}^{-1} = R_{-\theta} = \begin{pmatrix} \cos\theta & \sin\theta \\ -\sin\theta & \cos\theta \end{pmatrix}
\]

Applying $R_{\theta}^{-1}$ to $\vec{r}$ yields:
\[
  L_{R_{\theta}^{-1}}(\vec{r}) =  R_{\theta}^{-1} \begin{pmatrix} x \\ y \end{pmatrix} = \begin{pmatrix} x\cos\theta + y\sin\theta \\ -x\sin\theta + y\cos\theta \end{pmatrix}
\]

Thus, the rotated set $D_R$ is given by:
\begin{align*}
  D_R &= \left\{ \begin{pmatrix} x \\ y \end{pmatrix} \;\middle|\; f(x\cos\theta + y\sin\theta, -x\sin\theta + y\cos\theta) = 0 \right\} \\[1ex]
  \vec{u} &= \begin{pmatrix} x' \\ y' \end{pmatrix} = R_{\theta}^{-1}\vec{r} \iff \vec{r} = R_{\theta}\vec{u} \\[1ex]
  D_R &= \left\{ R_{\theta} \begin{pmatrix} x' \\ y' \end{pmatrix} \;\middle|\; f(x', y') = 0 \right\} = L_{R_{\theta}}(D)
\end{align*}

\end{document}
